\subsection{Criterios para que $\operatorname{Nor}(X)$ sea abierto.}
\label{IV.2.s6.13-fr}

\begin{env}[6.13.1]
\label{IV.2.s6.13.1-fr}
Dado un preesquema localmente noetheriano $X$, denotaremos por
$\operatorname{Nor}(X)$ el conjunto de los puntos $x\in X$ en los que $X$
es \emph{normal}; este conjunto contiene a $\operatorname{Reg}(X)$ y está contenido
en el conjunto (abierto) de los puntos $x$ tales que $\mathcal O_{X,x}$ sea
un dominio entero (es decir, el conjunto de los puntos $x$ en los que $X$
es reducido y que pertenecen a una sola componente irreducible de $X$).
\end{env}
